\documentclass[11pt,a4paper]{article}

\usepackage{fontspec}
\usepackage{polyglossia}
\setmainlanguage{russian}
\setotherlanguage{english}
\setmainfont{PT Serif}
\setsansfont{Arial}
\setmonofont{Menlo}
\newfontfamily\cyrillicfont{PT Serif}
\newfontfamily\cyrillicfontsf{Arial}
\newfontfamily\cyrillicfonttt{Menlo}

\usepackage[a4paper,margin=19mm,headheight=14pt]{geometry}
\usepackage{microtype}
\usepackage{amsmath,amssymb}
\usepackage{booktabs}
\usepackage{tabularx}
\usepackage{array}
\usepackage{enumitem}
\usepackage{xcolor}
\usepackage{colortbl}
\usepackage{fancyhdr}
\usepackage{titlesec}
\usepackage{hyperref}
\usepackage{xurl}

\definecolor{navy}{HTML}{183153}
\definecolor{blue}{HTML}{2B6CB0}
\definecolor{lightblue}{HTML}{EDF5FC}
\definecolor{lightgray}{HTML}{F3F4F6}
\definecolor{midgray}{HTML}{697386}
\definecolor{green}{HTML}{287A55}

\hypersetup{
  colorlinks=true,
  linkcolor=navy,
  urlcolor=blue,
  citecolor=green,
  pdfauthor={GMM research project},
  pdftitle={Экспериментальное исследование локализованных моментов для GMM},
  pdfsubject={Калибровка и устойчивость к равномерному загрязнению}
}

\pagestyle{fancy}
\fancyhf{}
\lhead{\sffamily\small Локализованные моменты для GMM}
\rhead{\sffamily\small Экспериментальное исследование}
\cfoot{\sffamily\small\thepage}
\renewcommand{\headrulewidth}{0.3pt}

\titleformat{\section}
  {\Large\sffamily\bfseries\color{navy}}{\thesection}{0.6em}{}
\titleformat{\subsection}
  {\large\sffamily\bfseries\color{blue}}{\thesubsection}{0.6em}{}
\titlespacing*{\section}{0pt}{1.30em}{0.55em}
\titlespacing*{\subsection}{0pt}{0.95em}{0.35em}

\setlist[itemize]{leftmargin=1.35em,itemsep=0.18em,topsep=0.25em}
\renewcommand{\arraystretch}{1.14}
\setlength{\tabcolsep}{4.2pt}
\setlength{\parindent}{0pt}
\setlength{\parskip}{0.34em}
\emergencystretch=2em

\newcolumntype{Y}{>{\raggedright\arraybackslash}X}
\newcolumntype{R}{>{\raggedleft\arraybackslash}X}
\newcommand{\code}[1]{{\footnotesize\nolinkurl{#1}}}
\newcommand{\best}[1]{\textbf{\color{green}#1}}
\newcommand{\E}{\mathbb{E}}
\newcommand{\KL}{\mathrm{KL}}
\newcommand{\ARI}{\mathrm{ARI}}
\newcommand{\Mzero}{M_0}
\newcommand{\Mone}{M_1}
\newcommand{\Mtwo}{M_2}

\begin{document}

\begin{center}
  {\sffamily\bfseries\LARGE
    Экспериментальное исследование локализованных моментов\
    для оценивания гауссовских смесей\par}
  \vspace{0.6em}
  {\large Калибровка метода и устойчивость к равномерному загрязнению\par}
  \vspace{0.4em}
  {\small Зафиксированная версия результатов: 13 августа 2026 г.}
\end{center}

\begin{abstract}
Мы исследуем оценивание сферических гауссовских смесей по локализованным
нулевым, первым и вторым моментам. Предварительные эксперименты показывают,
что на корректно специфицированных данных преимущество над EM ограничено
смесями с почти одинаковыми дисперсиями. Напротив, при замене
$2$--$10\%$ наблюдений равномерным фоновым шумом критерий только на вторых
моментах устойчиво превосходит EM по восстановлению незагрязнённой смеси и
кластеризации. В подтверждающем эксперименте с одинаковым набором тестовых
функций и новыми seed при $5\%$ загрязнения средний forward KL уменьшается с
$0.277$ до $0.145$, а ARI возрастает с $0.384$ до $0.534$. Контроли показывают,
что эффект не объясняется большим числом направлений, случайностью запуска
или размещением тестов в областях малой плотности. Добавление нулевых и первых
моментов при достаточно сильном загрязнении ухудшает результат; следовательно,
это не безусловно полезные дополнительные ограничения, а чувствительные к
misspecification компоненты критерия. Отдельно выводится leave-one-out
поправка, устраняющая смещение, возникающее, когда центрами тестов служат сами
наблюдения.
\end{abstract}

\section{Метод и экспериментальный протокол}

Для центра $c$, масштаба $s$ и единичных направлений $u,v$ использовались
нормированные по масштабу тестовые функции
\begin{align}
  \omega_0(x;c,s) &= K_s(x-c), \\
  \omega_1(x;c,s,u) &=
    \frac{\langle x-c,u\rangle}{s}K_s(x-c), \\
  \omega_2(x;c,s,u,v) &=
    \frac{\langle x-c,u\rangle\langle x-c,v\rangle}{s^2}K_s(x-c),
  \qquad
  K_s(z)=\exp\!\left(-\frac{\lVert z\rVert^2}{2s^2}\right).
\end{align}
Нормировка делает веса семейств моментов сопоставимыми при изменении $s$.
Для второго момента далее использовано $v=u$. Критерий имеет вид
\begin{equation}
  \mathcal L(\theta)=
  \sum_{r=0}^{2}\sum_{b=1}^{B}
  \frac{w_r}{2J_{rb}}
  \sum_{j=1}^{J_{rb}}\left(\widehat m_{rbj}-m_{rbj}(\theta)\right)^2,
  \label{eq:loss}
\end{equation}
где $b$ обозначает bandwidth-блок. Усреднение на $J_{rb}$ не позволяет числу
направлений неявно менять полный вес блока. Увеличение $J_{rb}$ всё же может
улучшать результат за счёт более
точного покрытия направлений и уменьшения Monte Carlo шума.

Основной сценарий содержит $K=3$ сбалансированных сферических компоненты в
$d=10$, $n=200$, фиксированную сложную геометрию средних с минимальным
махаланобисовым расстоянием $2$ и отношением дисперсий $3$. Доля
$\varepsilon\in\{0,0.02,0.05,0.10\}$ наблюдений заменяется равномерным шумом
в ограничивающем боксе с padding $4$. Все методы обучаются на полной
загрязнённой выборке, но KL и ошибки параметров вычисляются относительно
исходной гауссовской смеси, а ARI --- только по незагрязнённым наблюдениям.

Основная moment-модель использует все $200$ наблюдений как центры, два
масштаба, соответствующие $17$ и $133$ ожидаемым соседям, и по десять первых
и вторых направлений. Оптимизируются веса смеси, средние, сферические
дисперсии и масштабы; инициализация общая с baseline и основана на KMeans.
Сравниваются сферический EM с KMeans-инициализацией и
$\texttt{reg\_covar}\in\{10^{-6},10^{-2}\}$, а также KMeans. Во всех
таблицах меньшее значение KL и ошибок параметров лучше, большее ARI лучше.

Калибровочные эксперименты являются \emph{exploratory}. Главный эксперимент
на равномерном фоне использует новые $60$ dataset seed (101--160) и три
model seed на каждый датасет: всего $8640$ обучений для 12 моделей и четырёх
уровней загрязнения, все завершились успешно. Для парных сравнений сначала
усреднялись три model seed внутри датасета, затем строились percentile
bootstrap интервалы по 60 независимым датасетам. Интервалы номинально
95-процентные; поправка на множественные сравнения не применялась, поэтому
поиск лучшего веса внутри сетки следует считать описательным.

\section{Leave-one-out поправка для data-centered тестов}
\label{sec:loo}

Пусть центр $c_j=X_j$ совпадает с $j$-м наблюдением, а эмпирический отклик
вычисляется по той же выборке:
\begin{equation}
  \widehat m_r(X_j)=\frac1n\sum_{i=1}^{n}\omega_r(X_i;X_j),
  \qquad
  m_r(c)=\E_X\,\omega_r(X;c).
\end{equation}
Условно на $X_j$ центр уже не является независимым от слагаемых, и
\begin{equation}
  \E\!\left[\widehat m_r(X_j)\mid X_j\right]
  =\frac1n\omega_r(X_j;X_j)+\frac{n-1}{n}m_r(X_j).
  \label{eq:self-bias}
\end{equation}
Для нулевого момента $\omega_0(X_j;X_j)=1$, поэтому bias равен
$[1-m_0(X_j)]/n$ и обычно направлен вверх. Для первых и вторых моментов
self-term равен нулю, поскольку $X_j-c_j=0$, но проблема не исчезает:
деление на $n$ вместо $n-1$ даёт bias $-m_r(X_j)/n$. Следовательно, корректны
не только удаление self-term, но и перенормировка:
\begin{equation}
  \widehat m_r^{(-j)}(X_j)
  =\frac1{n-1}\sum_{i\ne j}\omega_r(X_i;X_j),
  \qquad
  \E\!\left[\widehat m_r^{(-j)}(X_j)\mid X_j\right]=m_r(X_j).
  \label{eq:loo}
\end{equation}
Эта же логика применена к оценке локального числа соседей при выборе $s$.
Для центров, сгенерированных независимо от данных, исключение наблюдения не
требуется. После введения поправки KL конфигурации $\Mzero+\Mone$ с широким
масштабом снизился с $0.100$ до $0.087$ при отношении дисперсий $1$ и с
$0.101$ до $0.088$ при отношении $1.2$. Это достаточно большой практический
эффект, поэтому все последующие результаты используют leave-one-out.

\section{Калибровка на корректно специфицированных данных}

\subsection{Порядки моментов и масштабы}

Сначала варьировались порядки моментов и локальность тестов на чистых
сферических смесях. В таблице~\ref{tab:clean-spread} приведён лучший
moment-вариант внутри exploratory grid и лучший EM. При почти равных
дисперсиях сочетание нулевого и трёх первых моментов на широком масштабе
превосходит EM. Эффект исчезает при сильном разбросе дисперсий.

\begin{table}[ht]
\centering
\small
\caption{Чистые данные: лучший moment-вариант в exploratory grid и EM.
Усреднение по 20 dataset seed и двум model seed.}
\label{tab:clean-spread}
\begin{tabularx}{\linewidth}{@{}lYrrrr@{}}
\toprule
Разброс $\sigma^2$ & Метод & KL & TV весов & Ошибка $\mu$ & ARI \\
\midrule
$1.0$ & $\Mzero+3\Mone$, $s(133)$ & \best{0.0867} & \best{0.0536} & \best{0.1846} & \best{0.3263} \\
      & EM, $10^{-2}$                 & 0.1076 & 0.1329 & 0.2245 & 0.2894 \\
\addlinespace
$1.2$ & $\Mzero+3\Mone$, $s(133)$ & \best{0.0883} & \best{0.0491} & \best{0.1812} & \best{0.3548} \\
      & EM, $10^{-2}$                 & 0.1079 & 0.1356 & 0.2166 & 0.2985 \\
\addlinespace
$6.0$ & лучший moment-вариант          & 0.1182 & 0.0526 & 0.1660 & 0.7402 \\
      & EM, $10^{-2}$                 & \best{0.0901} & \best{0.0487} & \best{0.1348} & \best{0.7698} \\
\bottomrule
\end{tabularx}
\end{table}

При промежуточном отношении дисперсий $3$ EM также остаётся лучше: его
KL/ARI равны $0.0869/0.6036$, тогда как лучший moment-вариант даёт примерно
$0.1052/0.5794$. Таким образом, преимущество на почти гомоскедастичных
чистых смесях реально, но слишком узко для основной мотивации метода.
Сочетание локального и широкого масштабов $[17,133]$ было лучше одного
масштаба и сетки $[17,67,133]$; это указывает на полезность двух разрешений,
а не просто на преимущество максимально малого bandwidth.

\subsection{Число вторых направлений}

При $w_0=0$, $w_1=1$, $w_2=8$ увеличение числа вторых направлений
$D_2=1,3,10,20$ изменяло KL как
$0.1133,0.1075,0.1058,0.1057$, а ARI как
$0.5745,0.5792,0.5794,0.5781$. Следовательно, основной выигрыш достигается
между одним и десятью направлениями, после чего наблюдается насыщение.
Десять направлений выбраны как практически достигнутая верхняя граница;
20 удваивают соответствующую часть вычислений без видимого улучшения.

\section{Равномерное загрязнение: основной результат}

\subsection{Равный бюджет moment-моделей}

В предварительном опыте вариант только со вторыми моментами оказался лучшим,
но строил больше активных тестов, чем варианты с первыми моментами. В
подтверждающем эксперименте все moment-модели поэтому \emph{конструируют один
и тот же} набор: один нулевой, десять первых и десять вторых тестов на каждом
центре и обоих масштабах. Неиспользуемое семейство зануляется только весом в
\eqref{eq:loss}. Для фиксированных dataset/model seed совпадают KMeans
инициализация, центры, масштабы и направления. При двух масштабах критерий
чистого $\Mtwo$ содержит $200\cdot10\cdot2=4000$ активных условий. Это
устраняет различия в числе проб и Monte Carlo направлениях внутри moment
абляции. Алгоритмический бюджет EM с этим числом не отождествляется: EM
остаётся самостоятельным и существенно более быстрым baseline.

\begin{table}[ht]
\centering
\small
\caption{Средние результаты на чистой смеси и при равномерном загрязнении.
Mixed означает $\Mtwo+w_1\Mone$ с репрезентативным малым $w_1$ из сетки;
выбор веса является описательным.}
\label{tab:main}
\begin{tabularx}{\linewidth}{@{}clrrrrr@{}}
\toprule
$\varepsilon$ & Метод & KL & TV весов & Ошибка $\mu$ & Ошибка $\Sigma$ & ARI \\
\midrule
0 & $\Mtwo$ & 0.1096 & 0.0659 & 0.1645 & 0.0878 & 0.5732 \\
  & Mixed, $w_1=1/4$ & 0.1049 & \best{0.0590} & 0.1590 & 0.0842 & \best{0.5795} \\
  & EM, $10^{-2}$ & \best{0.0989} & 0.0746 & \best{0.1564} & \best{0.0825} & 0.5729 \\
  & KMeans & 0.2290 & 0.1097 & 0.2239 & 0.2092 & 0.4865 \\
\addlinespace
.02 & $\Mtwo$ & 0.1155 & 0.0697 & 0.1698 & 0.1111 & 0.5657 \\
  & Mixed, $w_1=1/8$ & \best{0.1132} & \best{0.0694} & \best{0.1682} & \best{0.1107} & \best{0.5667} \\
  & EM, $10^{-2}$ & 0.1919 & 0.1969 & 0.4295 & 0.5373 & 0.4699 \\
  & KMeans & 0.2874 & 0.1151 & 0.2498 & 0.3065 & 0.4459 \\
\addlinespace
.05 & $\Mtwo$ & \best{0.1449} & \best{0.0873} & \best{0.2091} & \best{0.1889} & \best{0.5339} \\
  & Mixed, $w_1=1/32$ & 0.1487 & 0.0898 & 0.2102 & 0.1992 & 0.5301 \\
  & EM, $10^{-2}$ & 0.2775 & 0.2917 & 0.5019 & 1.0959 & 0.3845 \\
  & KMeans & 0.4407 & 0.1314 & 0.3108 & 0.4682 & 0.3871 \\
\addlinespace
.10 & $\Mtwo$ & \best{0.3208} & \best{0.1458} & \best{0.3300} & \best{0.5305} & \best{0.4023} \\
  & Mixed, $w_1=1/32$ & 0.3367 & 0.1528 & 0.3282 & 0.5574 & 0.3957 \\
  & EM, $10^{-2}$ & 0.3387 & 0.2540 & 0.4286 & 1.1268 & 0.3688 \\
  & KMeans & 0.7712 & 0.1709 & 0.5232 & 0.7532 & 0.2888 \\
\bottomrule
\end{tabularx}
\end{table}

На чистых данных EM имеет небольшое преимущество по KL, но методы близки по
ARI. Уже при $2\%$ загрязнения $\Mtwo$ лучше EM по всем пяти внешним
метрикам; разрыв максимален при $5\%$. При $10\%$ преимущество по KL
сокращается до $0.018$, но сохраняется по параметрам и ARI. Парные разности
$\Mtwo-\mathrm{EM}$ с bootstrap интервалами приведены в
таблице~\ref{tab:paired}. Для $\varepsilon\geq0.02$ интервалы не пересекают
нуль по всем пяти метрикам. Повторение предварительного среднего KL чистого
$\Mtwo$ на новых seed с отклонением менее $0.004$ дополнительно исключает
объяснение результата единичной удачной выборкой.

\begin{table}[ht]
\centering
\small
\caption{Парная разность $\Mtwo-\mathrm{EM}_{10^{-2}}$; в скобках ---
номинальный 95\% bootstrap CI по 60 датасетам. Отрицательная разность KL
и положительная разность ARI означают преимущество $\Mtwo$.}
\label{tab:paired}
\begin{tabularx}{0.94\linewidth}{@{}cRR@{}}
\toprule
$\varepsilon$ & $\Delta$ KL & $\Delta$ ARI \\
\midrule
0    & $+0.0107\ [0.0041,\,0.0170]$ & $+0.0003\ [-0.0098,\,0.0104]$ \\
.02  & $-0.0764\ [-0.0898,\,-0.0630]$ & $+0.0959\ [0.0728,\,0.1189]$ \\
.05  & $-0.1326\ [-0.1437,\,-0.1213]$ & $+0.1494\ [0.1277,\,0.1715]$ \\
.10  & $-0.0179\ [-0.0345,\,-0.0015]$ & $+0.0334\ [0.0098,\,0.0575]$ \\
\bottomrule
\end{tabularx}
\end{table}

\subsection{Почему дополнительные порядки моментов могут вредить}

На чистых данных небольшой первый момент полезен: при $w_1=1/4$ он снижает
KL относительно чистого $\Mtwo$ на $0.0047$ с CI
$[-0.0070,-0.0029]$ и повышает ARI на $0.0063$ с CI
$[0.0020,0.0108]$. При $2\%$ оптимальный положительный вес уменьшается примерно
до области $1/8$--$1/4$. При $5\%$ даже минимальный проверенный $w_1=1/32$
уже увеличивает
KL на $0.0038$ и уменьшает ARI на $0.0038$; при $10\%$ соответствующие
изменения равны $+0.0160$ и $-0.0066$, и интервалы не содержат нуль.
Эмпирически оптимальный путь поэтому имеет вид
\begin{equation}
  w_1/w_2:\qquad 1/4\ \longrightarrow\ [1/8,1/4]\ \longrightarrow\ 0
  \quad\text{при росте }\varepsilon.
\end{equation}

Нулевой момент ещё чувствительнее. При $w_0=1/16$, $w_1=0$, $w_2=1$ KL/ARI
равны $0.1082/0.5740$ на чистых данных. При загрязнении
$\varepsilon=0.02,0.05,0.10$ они становятся соответственно
$0.1430/0.5275$, $0.2861/0.4021$ и $0.4258/0.3176$. Уже при $2\%$ добавление
этого момента к $\Mtwo$ увеличивает KL на $0.0275$ и уменьшает ARI на
$0.0382$ с интервалами, не
пересекающими нуль. Чистый первый момент также хуже чистого второго на всех
уровнях: например, при $5\%$ его KL/ARI равны $0.2278/0.4684$.

Этот эффект не противоречит тому, что разные моменты несут дополнительную
информацию при корректной модели. При misspecification эмпирические отклики
соответствуют загрязнённому распределению, тогда как параметрическое семейство
содержит только три гауссовские компоненты. Все уравнения одновременно
выполнить нельзя, и добавление семейства меняет pseudo-true минимум
\eqref{eq:loss}. Нулевой момент непосредственно измеряет локальную массу и
поэтому особенно реагирует на добавленный фон. Первые моменты чувствительны к
локальному сдвигу массы и границам бокса. Локализованный второй момент в
исследованной геометрии лучше сохраняет информацию о масштабе и расположении
ядра смеси. Последнее является эмпирическим объяснением, а не доказанной
универсальной робастностью: функция $t^2e^{-t^2/2}$ ограничена, но критерий не
является формальным robust estimator с заданной influence function.

\section{Откуда берётся устойчивость: контроль центров}

Возникла гипотеза, что преимущество $\Mtwo$ связано не с порядком момента, а
с тестами, случайно попавшими в области малой плотности: почти нулевой отклик
мог бы сообщать модели, где плотности быть не должно. Для проверки при
$\varepsilon=0.05$ сравнивались по 200 пространственных центров:
\begin{itemize}
  \item все наблюдаемые точки (190 чистых и 10 выбросов);
  \item только oracle-чистые центры с повторением десяти из них;
  \item 190 чистых центров и 10 независимых равномерных точек из того же бокса;
  \item точная контрольная реконструкция наблюдаемых центров.
\end{itemize}
Дополнительно на чистых данных сравнивались обычные центры и 10 искусственных
равномерных центров. Эксперимент использовал 40 новых dataset seed (201--240),
три model seed и завершил все 3600 обучений.

Малый-масштабный отклик $\Mtwo$ действительно различает области: его среднее
равно $0.0310$ на чистых центрах, $2.78\cdot10^{-4}$ на наблюдаемых выбросах
и $9.25\cdot10^{-5}$ на независимых равномерных центрах. Однако такие тесты
не улучшают оценивание. Относительно oracle-чистых центров добавление
независимых равномерных центров увеличивает KL на $0.0068$ и уменьшает ARI на
$0.0064$; интервалы не содержат нуль. На полностью чистых данных их эффект на
KL и ARI статистически неразличим с нулём, а TV и ошибка средних слегка
ухудшаются. Точная контрольная реконструкция совпала со стандартным режимом,
что также проверяет реализацию эксперимента.

Удаление центров в наблюдаемых выбросах, напротив, немного помогает. Чтобы
отделить этот эффект от изменения выбранных bandwidth, был проведён прямой
контроль с фиксированными масштабами (720 дополнительных обучений):

\begin{table}[ht]
\centering
\small
\caption{Контроль центров при $5\%$ загрязнения. Вторая строка меняет только
центры, сохраняя масштабы стандартного режима; третья также перевыбирает
масштабы по чистым центрам.}
\label{tab:centers}
\begin{tabularx}{\linewidth}{@{}Yrrrrr@{}}
\toprule
Центры и масштабы & KL & TV весов & Ошибка $\mu$ & Ошибка $\Sigma$ & ARI \\
\midrule
Все наблюдения, observed $s$ & 0.1496 & 0.0883 & 0.1871 & 0.1696 & 0.5528 \\
Oracle clean, тот же $s$ & \best{0.1445} & \best{0.0843} & \best{0.1836} & 0.1587 & 0.5573 \\
Oracle clean, собственный $s$ & 0.1438 & 0.0850 & 0.1860 & \best{0.1578} & \best{0.5577} \\
\bottomrule
\end{tabularx}
\end{table}

При неизменных $s$ замена outlier-центров на чистые снижает KL на $0.00511$;
границы 95\% CI равны $-0.00660$ и $-0.00389$. Ошибка ковариаций снижается
на $0.01086$,
а ARI повышается на
$0.00446$; KL улучшается на 39 из 40 датасетов. Повторный выбор bandwidth
даёт лишь $-0.00063$ KL с интервалом, содержащим нуль. Значит, outlier-центры
скорее немного вредят, а «информация о пустом пространстве» не объясняет
основной выигрыш над EM. Oracle-фильтрация улучшает KL только примерно на
$3.4\%$ и не может объяснить разрыв $0.145$ против $0.277$. Более того,
положительный вес первого момента остаётся вредным и с clean-only центрами:
его проблема вызвана откликом загрязнения и в тестах вокруг чистых точек, а не
только центрами на выбросах.

\section{Связь с более ранними экспериментами}

Широкие ранние sweeps не показывали преимущества над EM на равномерном фоне.
Это не противоречит основному результату: там одновременно использовались все
порядки моментов с равными весами, автоматический преимущественно
средний/широкий bandwidth, мало вторых направлений, короткая оптимизация и
оцениватель без полной leave-one-out коррекции. Кроме того, старые uniform
сценарии меняли вместе с $\varepsilon$ размер выборки, размерность и padding,
поэтому не образуют сопоставимой кривой загрязнения. Новый результат появился
не из oracle-центров: стандартный, полностью наблюдаемый $\Mtwo$ уже
превосходит EM. Наиболее существенное изменение --- удаление чувствительных
$\Mzero$ и $\Mone$, после чего результат подтверждён на новых seed и при
равном числе построенных проб.

Два других misspecified семейства дали более смешанную картину. На
Student-$t$ компонентах с тремя степенями свободы лучший наблюдавшийся
moment-ARI был около $0.533$ против $0.342$ у EM, но даже лучший moment-KL
равнялся примерно $1.508$ против $1.002$. На направленно асимметричных
Gamma-возмущениях лучший moment-вариант не превзошёл EM: средний true KL был
примерно $0.283$ против $0.242$. Фильтрация центров по локальной плотности
иногда стабилизировала тяжёлые хвосты, но не давала универсального улучшения.
Эти отрицательные результаты показывают, что наблюдаемая устойчивость пока
является конкретной для diffuse gross contamination, а не общей победой при
любом отклонении от гауссовости.

\section{Малый training loss и риск переобучения}

В основном опыте при $5\%$ загрязнения финальное значение критерия чистого
$\Mtwo$ равно в среднем $3.93\cdot10^{-6}$ (в отдельном контроле центров
$3.78\cdot10^{-6}$). Это число выглядит необычно малым, но его масштаб
обусловлен нормировкой: при двух равновесных блоках вторых моментов
\eqref{eq:loss} является средним квадратом отклика по 4000 условиям. Поэтому
$\sqrt{\mathcal L}\approx1.9\cdot10^{-3}$, тогда как характерный эмпирический
отклик малого масштаба около $3.1\cdot10^{-2}$. Кроме того, оптимизация
уменьшает начальный loss лишь примерно вдвое, а не на несколько порядков.

Модель имеет около 35 свободных параметров против 4000 условий, хотя сами
условия сильно коррелированы. Малый training loss поэтому не доказывает ни
интерполяцию, ни хорошее восстановление распределения. Действительно, среди
120 запусков чистого $\Mtwo$ корреляция final loss с KL составляет лишь около
$0.22$, с ошибкой ковариаций $0.13$, а с ARI $0.03$. Leave-one-out устраняет
self-bias из раздела~\ref{sec:loo}, но не устраняет оптимизм от выбора
параметров и тестов на одной выборке.

В то же время выигрыш воспроизводится на независимых сгенерированных
датасетах, поэтому он не сводится к запоминанию одной фиксированной выборки.
Корректный следующий контроль --- разделять данные на fit/test части или
оценивать моментные невязки на независимой большой выборке при фиксированных
центрах и масштабах. До такого эксперимента final loss следует трактовать
только как диагностический критерий оптимизации, а внешние KL, ошибки
параметров и ARI --- как основные показатели качества.

\section{Выводы и ограничения}

Эксперименты поддерживают следующий ограниченный, но содержательный вывод.
Локализованные вторые моменты дают практически значимое преимущество над EM
при небольшом числе наблюдений и $2$--$10\%$ диффузного равномерного
загрязнения, особенно при $5\%$. Эффект воспроизводится на новых seed,
сохраняется при одинаковом вычислительном бюджете moment-вариантов и не
объясняется искусственными тестами в пустом пространстве. На чистых данных EM
остаётся немного лучше в сценарии с неодинаковыми дисперсиями; слабый первый
момент полезен только при малом загрязнении, а нулевой быстро становится
вредным.

Для статьи этот результат лучше формулировать как преимущество в конкретном
режиме misspecification, а не как общее превосходство над maximum likelihood.
Главные ограничения текущей серии: один основной размер
$(d,n,K)=(10,200,3)$, один тип и геометрия загрязнения, spherical fit,
post-hoc выбор веса первого момента и отсутствие held-out moment loss.
Приоритетные проверки перед финальными экспериментами --- независимая
валидация moment-критерия, перенос фиксированной конфигурации чистого $\Mtwo$
на разные $d,n,K$, padding и структуры смеси, а затем тяжёлые хвосты без
дополнительной настройки.

\section*{Воспроизводимость}

{\small
Конфигурации находятся в \code{experiments/configs/}, результаты --- в
одноимённых каталогах \code{results/}. Основные run ID:
\code{moments_uniform_background_equal_budget_confirmation},
\code{moments_negative_space_center_mechanism} и
\code{moments_negative_space_fixed_bandwidth_control}. Для калибровки
использованы \code{moments_order_bandwidth_ablation_loo} и три run ID с
префиксом \code{moments_varspread3_}: \code{bandwidth_loo},
\code{order_weights_loo}, \code{second_weight_directions_loo}.
}

\end{document}
